\chapter{Mathematical induction}
\chaptercontents{A & Weak induction (\S4.2) \\ B & Strong induction, well-ordering (\S4.3) \\ C & Chapter 4 exercises (4.E)}%
{\fE\ 4.2: 5, 11 \textperiodcentered\ 4.3: 8, 9, 14, 15 \textperiodcentered\ 4.E: 10, 11, 13, 16--20, 25--27, 29--32\\ \fR\ 4.2: 8, 17, 20 \textperiodcentered\ 4.3: 10, 18 \textperiodcentered\ 4.E: 12, 14, 15, 21, 24, 28, 33--35\\ \S4.1 is not in the planning. Tested in the \textbf{exam} (4 Nov).}

\hsection{Weak induction}

\begin{key}[{What you quote in this section}]
\textbf{Weak induction} (Theorem 4.2.1, Strategy 4.2.2): to prove $\forall n\ge n_0,\ p(n)$, prove the \emph{base case} $p(n_0)$ and the \emph{induction step}: for all $n\ge n_0$, $p(n)\Rightarrow p(n+1)$. Write ``We proceed by induction on $n\ge n_0$'', state $p(n)$, then ``(Base case)'', ``(Induction step) Let $n\ge n_0$ and suppose $p(n)$ \dots\ we must prove $p(n+1)$ \dots'', then ``the result follows by induction''. \textbf{Recursive sums} (Definition 4.1.10): $\sum_{k=0}^{n+1}a_k=\big(\sum_{k=0}^na_k\big)+a_{n+1}$. \textbf{Binomial coefficients} (Definition 4.1.15, Theorem 4.2.15): $\binom n0=1$, $\binom0{k+1}=0$, $\binom{n+1}{k+1}=\binom nk+\binom n{k+1}$; $\binom nk=\frac{n!}{k!(n-k)!}$ for $k\le n$ and $0$ for $k>n$. \textbf{Binomial theorem} (4.2.18): $(x+y)^n=\sum_{k=0}^n\binom nkx^ky^{n-k}$.
\end{key}

\begin{bookex}[essential]{4.2.5}
Prove by induction that $\displaystyle\sum_{k=0}^n2^k=2^{n+1}-1$ for all $n\in\N$.
\solution
\given Definition 4.1.10 for the recursive meaning of $\sum$.
\goal $\forall n\in\N,\ p(n)$ where $p(n)$: $\sum_{k=0}^n2^k=2^{n+1}-1$.
\begin{steps}
\item We proceed by induction on $n\ge0$, with $p(n)$ as above.
\item \emph{(Base case $n=0$.)} $\sum_{k=0}^02^k=2^0=1$, and $2^{0+1}-1=2-1=1$. So $p(0)$ holds.
\item \emph{(Induction step.)} Let $n\ge0$ and suppose $p(n)$: $\sum_{k=0}^n2^k=2^{n+1}-1$ \why{induction hypothesis}. We must prove $p(n+1)$: $\sum_{k=0}^{n+1}2^k=2^{n+2}-1$.
\item
\begin{align*}
\sum_{k=0}^{n+1}2^k&=\Big(\sum_{k=0}^n2^k\Big)+2^{n+1}\why{Definition 4.1.10}\\
&=(2^{n+1}-1)+2^{n+1}\why{induction hypothesis}\\
&=2\cdot2^{n+1}-1=2^{n+2}-1\why{algebra}.
\end{align*}
So $p(n+1)$ holds.
\end{steps}
\conclusion By induction, $\sum_{k=0}^n2^k=2^{n+1}-1$ for all $n\in\N$.\qed
\begin{compare}
\item The four sentences: ``by induction on $n$'' with $p(n)$ stated; base case computed on both sides; step with the hypothesis \emph{marked} where it is used; closing sentence.
\item Never start the step from the goal $\sum_{k=0}^{n+1}2^k=2^{n+2}-1$ and manipulate it; start from the left side and arrive at the right side.
\end{compare}
\end{bookex}

\begin{bookex}[recommended]{4.2.8}
Write down the statement $p(n)$ that Example 4.2.7 (``all horses have the same colour'') tries to prove for all $n\ge1$. Write down, with quantifiers, what the induction step must prove. Explain why the argument fails.
\solution
\begin{steps}
\item $p(n)$: ``every set of $n$ horses has the same colour'' (all horses in any $n$-element set share one colour).
\item \emph{Base case} $p(1)$: a set with one horse has only one colour. Correct.
\item \emph{Induction step} must prove $\forall n\ge1,\ \big(p(n)\Rightarrow p(n+1)\big)$: for \emph{every} $n\ge1$, if every $n$-set is monochrome then every $(n+1)$-set is.
\item \emph{The argument:} given $n+1$ horses $h_1,\dots,h_{n+1}$, the sets $\{h_1,\dots,h_n\}$ and $\{h_2,\dots,h_{n+1}\}$ are monochrome by $p(n)$; they overlap, so the colour of the overlap forces all $n+1$ to agree.
\item \emph{Why it fails:} for $n=1$ the two sets are $\{h_1\}$ and $\{h_2\}$, which do \emph{not} overlap, so nothing links the colour of $h_1$ to that of $h_2$. The implication $p(1)\Rightarrow p(2)$ is never proved, and the step only works for $n\ge2$. With the first link missing, induction cannot start: $p(2)$ is in fact false.
\end{steps}
\conclusion The step must hold for every $n\ge n_0$ including $n=n_0$; here it fails at $n=1$.\qed
\end{bookex}

\begin{bookex}[essential]{4.2.11}
Let $a,r\in\R$ with $r\ne1$. Prove that $\displaystyle\sum_{k=0}^nar^k=\frac{a(1-r^{n+1})}{1-r}$ for all $n\in\N$.
\solution
\given $r\ne1$, so $1-r\ne0$ and division by $1-r$ is allowed.
\begin{steps}
\item We proceed by induction on $n\ge0$ with $p(n)$: $\sum_{k=0}^nar^k=\frac{a(1-r^{n+1})}{1-r}$.
\item \emph{(Base case.)} $\sum_{k=0}^0ar^k=ar^0=a$, and $\frac{a(1-r^1)}{1-r}=\frac{a(1-r)}{1-r}=a$ \why{$1-r\ne0$}. So $p(0)$ holds.
\item \emph{(Induction step.)} Let $n\ge0$ and suppose $p(n)$. We must prove $\sum_{k=0}^{n+1}ar^k=\frac{a(1-r^{n+2})}{1-r}$.
\item
\begin{align*}
\sum_{k=0}^{n+1}ar^k&=\sum_{k=0}^nar^k+ar^{n+1}\why{Definition 4.1.10}\\
&=\frac{a(1-r^{n+1})}{1-r}+ar^{n+1}\why{induction hypothesis}\\
&=\frac{a(1-r^{n+1})+ar^{n+1}(1-r)}{1-r}\why{common denominator}\\
&=\frac{a-ar^{n+1}+ar^{n+1}-ar^{n+2}}{1-r}=\frac{a(1-r^{n+2})}{1-r}.
\end{align*}
So $p(n+1)$ holds.
\end{steps}
\conclusion By induction the formula holds for all $n\in\N$. (Exercise 4.2.5 is the case $a=1$, $r=2$.)\qed
\begin{compare}
\item Say where $r\ne1$ is used (base case and every division by $1-r$).
\end{compare}
\end{bookex}

\begin{bookex}[recommended]{4.2.17}
Prove that $\binom nk=\binom n{n-k}$ for all $n,k\in\N$ with $k\le n$.
\solution
\given Theorem 4.2.15: $\binom nk=\frac{n!}{k!(n-k)!}$ whenever $k\le n$.
\begin{steps}
\item Let $n,k\in\N$ with $k\le n$. Then $n-k\in\N$ and $n-k\le n$, so Theorem 4.2.15 applies to both $\binom nk$ and $\binom n{n-k}$.
\item $\binom n{n-k}=\dfrac{n!}{(n-k)!\,(n-(n-k))!}=\dfrac{n!}{(n-k)!\,k!}=\dfrac{n!}{k!\,(n-k)!}=\binom nk$ \why{$n-(n-k)=k$; commutativity of multiplication}.
\end{steps}
\conclusion $\binom nk=\binom n{n-k}$.\qed
\begin{compare}
\item No induction is needed once Theorem 4.2.15 is available; the check $n-k\le n$ makes the theorem applicable.
\end{compare}
\end{bookex}

\begin{bookex}[recommended]{4.2.20}
Give a proof of $\displaystyle\sum_{i=0}^n(-1)^i\binom ni=0$ using the binomial theorem.
\solution
\given Theorem 4.2.18 with $x=-1$, $y=1$. The identity is meant for $n\ge1$ (for $n=0$ the sum is $1$).
\begin{steps}
\item Let $n\ge1$. By the binomial theorem with $x=-1$, $y=1$:
\[
(-1+1)^n=\sum_{i=0}^n\binom ni(-1)^i1^{n-i}=\sum_{i=0}^n(-1)^i\binom ni \why{$1^{n-i}=1$}.
\]
\item The left side is $0^n=0$ because $n\ge1$.
\end{steps}
\conclusion $\sum_{i=0}^n(-1)^i\binom ni=0$ for all $n\ge1$. (Interpretation: a finite non-empty set has as many even-sized as odd-sized subsets.)\qed
\begin{compare}
\item State the substitution ($x=-1$, $y=1$) explicitly and note why $n\ge1$ is needed.
\end{compare}
\end{bookex}

\hsection{Strong induction, well-ordering}

\begin{key}[{What you quote in this section}]
\textbf{Strong induction} (Theorem 4.3.2, Strategies 4.3.3 and 4.3.6): to prove $\forall n\ge n_0,\ p(n)$, prove base cases $p(n_0),\dots,p(n_1)$ and the step: ``fix $n\ge n_1$ and assume $p(k)$ for all $n_0\le k\le n$; prove $p(n+1)$''. Use it when $p(n+1)$ depends on $p(n-1)$, $p(n-2)$, \dots; then you need as many base cases as the recursion looks back, and the step starts at $n_1$ so that the recursion applies to $n+1$. \textbf{Well-ordering} (Theorem 4.3.11): every non-empty subset of $\N$ has a least element. \textbf{Infinite descent} (Definition 4.3.16): $X\subseteq\R$ has infinite descent if $\forall x\in X,\ \exists y\in X,\ y<x$; no non-empty subset of $\N$ has it.
\end{key}

\begin{bookex}[essential]{4.3.8}
Define $a_0=4$, $a_1=9$ and $a_n=5a_{n-1}-6a_{n-2}$ for $n\ge2$. Prove that $a_n=3\cdot2^n+3^n$ for all $n\in\N$.
\solution
\given The recursion uses two previous terms, so two base cases and strong induction with the step starting at $n=1$ (Example 4.3.7 is the model).
\begin{steps}
\item We proceed by strong induction on $n\ge0$ with $p(n)$: $a_n=3\cdot2^n+3^n$, base cases $n=0,1$.
\item \emph{(Base cases.)} $a_0=4=3\cdot1+1=3\cdot2^0+3^0$. $a_1=9=6+3=3\cdot2^1+3^1$.
\item \emph{(Induction step.)} Fix $n\ge1$ and assume $a_k=3\cdot2^k+3^k$ for all $0\le k\le n$. We must prove $a_{n+1}=3\cdot2^{n+1}+3^{n+1}$. Since $n+1\ge2$ the recursion applies:
\begin{align*}
a_{n+1}&=5a_n-6a_{n-1}\why{definition, valid as $n+1\ge2$}\\
&=5(3\cdot2^n+3^n)-6(3\cdot2^{n-1}+3^{n-1})\why{hypothesis for $k=n$ and $k=n-1$; both $\ge0$}\\
&=15\cdot2^n+5\cdot3^n-9\cdot2^n-2\cdot3^n\why{$6\cdot2^{n-1}=3\cdot2^n$, $6\cdot3^{n-1}=2\cdot3^n$}\\
&=6\cdot2^n+3\cdot3^n=3\cdot2^{n+1}+3^{n+1}.
\end{align*}
\end{steps}
\conclusion By strong induction, $a_n=3\cdot2^n+3^n$ for all $n\in\N$.\qed
\begin{compare}
\item Two base cases, step fixed at $n\ge1$ (so that $n-1\ge0$ and $n+1\ge2$), and the hypothesis used for \emph{two} values of $k$.
\end{compare}
\end{bookex}

\begin{bookex}[essential]{4.3.9}
The Tribonacci sequence is defined by $t_0=0$, $t_1=0$, $t_2=1$ and $t_n=t_{n-1}+t_{n-2}+t_{n-3}$ for $n\ge3$. Prove that $t_n\le2^{n-3}$ for all $n\ge3$.
\solution
\given The recursion looks back three terms, so three base cases $n=3,4,5$, and the step is fixed at $n\ge5$ (so that $n-2\ge3$ and the hypothesis applies to all three terms).
\begin{steps}
\item We proceed by strong induction on $n\ge3$ with $p(n)$: $t_n\le2^{n-3}$, base cases $n=3,4,5$.
\item \emph{(Base cases.)} $t_3=t_2+t_1+t_0=1\le1=2^0$; $t_4=t_3+t_2+t_1=2\le2=2^1$; $t_5=t_4+t_3+t_2=4\le4=2^2$.
\item \emph{(Induction step.)} Fix $n\ge5$ and assume $t_k\le2^{k-3}$ for all $3\le k\le n$. We must prove $t_{n+1}\le2^{n-2}$.
\item Since $n+1\ge6\ge3$, $t_{n+1}=t_n+t_{n-1}+t_{n-2}$ \why{definition}. The indices $n,n-1,n-2$ are all $\ge3$ \why{$n\ge5$} and $\le n$, so the hypothesis applies to each:
\[
t_{n+1}\le2^{n-3}+2^{n-4}+2^{n-5}=2^{n-5}(4+2+1)=7\cdot2^{n-5}<8\cdot2^{n-5}=2^{n-2}.
\]
\end{steps}
\conclusion By strong induction, $t_n\le2^{n-3}$ for all $n\ge3$.\qed
\begin{compare}
\item Explain why three base cases: the step at $n=5$ uses $t_5,t_4,t_3$; a step at $n=4$ would need $t_2$, for which $p$ is not stated.
\end{compare}
\end{bookex}

\begin{bookex}[recommended]{4.3.10}
Prove that for all $n\ge8$, a value of $n$ dubloons can be obtained using only $3$-dubloon and $5$-dubloon coins.
\solution
\given $p(n)$: $\exists a,b\in\N,\ n=3a+5b$. The hint: $p(n)\Rightarrow p(n+3)$ (add one $3$-coin). So the step for $n+1$ uses $p(n-2)$: strong induction with base cases $8,9,10$ and the step fixed at $n\ge10$.
\begin{steps}
\item We proceed by strong induction on $n\ge8$ with base cases $8,9,10$.
\item \emph{(Base cases.)} $8=3+5$ ($a=b=1$); $9=3\cdot3$ ($a=3$, $b=0$); $10=5\cdot2$ ($a=0$, $b=2$).
\item \emph{(Induction step.)} Fix $n\ge10$ and assume $p(k)$ for all $8\le k\le n$. We must prove $p(n+1)$.
\item Since $n\ge10$, $n-2\ge8$, so $p(n-2)$ is available: $n-2=3a+5b$ for some $a,b\in\N$ \why{Strategy 1.2.30}.
\item Then $n+1=(n-2)+3=3(a+1)+5b$ with $a+1,b\in\N$, so $p(n+1)$ holds \why{$\exists$I}.
\end{steps}
\conclusion By strong induction, every $n\ge8$ is a sum of $3$s and $5$s. (The statement is false for $7$: $7=3a+5b$ would need $b\le1$, giving $3a=7$ or $3a=2$.)\qed
\begin{compare}
\item Three base cases because the step jumps back by $3$; the step must be fixed at $n\ge10$ so that $n-2\ge8$.
\end{compare}
\end{bookex}

\begin{bookex}[essential]{4.3.14}
What goes wrong in the proof of Proposition 4.3.12 ($\sqrt2$ is irrational) if we try to prove that $\sqrt4$ is irrational?
\solution
\given The proof of 4.3.12: write $\sqrt2=\frac ab$ in lowest terms (Lemma 4.3.13); $a^2=2b^2$; so $a^2$ is even, so $a$ is even (Proposition 1.1.63), $a=2c$; then $4c^2=2b^2$, $b^2=2c^2$, so $b$ is even; $2$ divides both: contradiction.
\begin{steps}
\item Run the same argument: suppose $\sqrt4=\frac ab$ in lowest terms. Squaring: $a^2=4b^2$. Then $a^2=2(2b^2)$ is even, so $a$ is even \why{Proposition 1.1.63}: $a=2c$.
\item Substitute: $4c^2=4b^2$, so $c^2=b^2$.
\item In the $\sqrt2$ proof the corresponding equation was $b^2=2c^2$, which forced $b$ even. Here $c^2=b^2$ says nothing about the parity of $b$. No contradiction can be derived; the argument simply stops.
\item And indeed there is none to derive: $\sqrt4=2=\frac21$ with $a=2$, $b=1$ in lowest terms satisfies everything ($c=1$, $c^2=b^2$).
\end{steps}
\conclusion The step ``hence $b$ is even'' fails, because after cancelling, the factor $2$ no longer appears on the right-hand side (it appeared an even number of times in $4=2^2$).\qed
\end{bookex}

\begin{bookex}[essential]{4.3.15}
Prove that $\sqrt3$ is irrational.
\solution
\given Lemma 4.3.13 (a positive rational is $\frac ab$ with $a,b\in\N$ non-zero and $1$ their only common natural divisor). We need the analogue of Proposition 1.1.63 for $3$, proved as a lemma first (the book's writing tip).
\begin{steps}
\item \emph{Lemma: for $a\in\Z$, if $3\mid a^2$ then $3\mid a$.} By Theorem 0.42, $a=3q+r$ with $r\in\{0,1,2\}$. If $r=1$: $a^2=9q^2+6q+1=3(3q^2+2q)+1$, remainder $1$. If $r=2$: $a^2=9q^2+12q+4=3(3q^2+4q+1)+1$, remainder $1$. In both cases $3\nmid a^2$ \why{the remainder is $1\ne0$; uniqueness in Theorem 0.42}. So $3\mid a^2$ forces $r=0$, i.e.\ $3\mid a$ \why{contraposition}.
\item \emph{Main proof.} Suppose, towards a contradiction, that $\sqrt3\in\Q$ \why{Strategy 1.1.51}. Since $\sqrt3>0$, Lemma 4.3.13 gives $\sqrt3=\frac ab$ with $a,b\in\N\setminus\{0\}$ having no common natural divisor other than $1$.
\item Squaring: $3=\frac{a^2}{b^2}$, so $a^2=3b^2$. Hence $3\mid a^2$ \why{Definition 0.36}, so $3\mid a$ \why{Lemma}: $a=3c$ with $c\in\Z$.
\item Substitute: $9c^2=3b^2$, so $b^2=3c^2$. Hence $3\mid b^2$, so $3\mid b$ \why{Lemma}.
\item Now $3$ is a natural number dividing both $a$ and $b$, and $3\ne1$: this contradicts the choice of $a,b$ in Step 2 \why{$\neg$E}.
\end{steps}
\conclusion $\sqrt3$ is irrational \why{$\neg$I}.\qed
\begin{compare}
\item The lemma ``$3\mid a^2\Rightarrow3\mid a$'' must be proved (by remainders), not assumed; it is used \emph{twice}.
\item Say where lowest terms (Lemma 4.3.13) is used: in the final contradiction.
\end{compare}
\end{bookex}

\begin{bookex}[recommended]{4.3.18}
Prove that the interval $(0,1)$ has infinite descent.
\solution
\given Definition 4.3.16: $\forall x\in(0,1),\ \exists y\in(0,1),\ y<x$.
\begin{steps}
\item Let $x\in(0,1)$ \why{Strategy 1.2.10}. Define $y=\frac x2$.
\item $y>0$ \why{$x>0$}; $y<x$ \why{$\frac x2<x$ as $x>0$}; $y<x<1$. So $y\in(0,1)$ and $y<x$ \why{Strategy 1.2.24}.
\end{steps}
\conclusion $(0,1)$ has infinite descent. (No contradiction arises, since $(0,1)\nsubseteq\N$; the well-ordering principle only forbids infinite descent in non-empty subsets of $\N$.)\qed
\end{bookex}

\hsection{Chapter 4 exercises}

\begin{bookex}[essential]{4.E.10}
Find all natural numbers $n$ such that $n^5<5^n$.
\solution
\plan Test small $n$ by hand to find where the inequality holds; then prove by induction that it holds from some point on.
\begin{steps}
\item Small cases: $n=0$: $0<1$ \tick; $n=1$: $1<5$ \tick; $n=2$: $32<25$ \cross; $n=3$: $243<125$ \cross; $n=4$: $1024<625$ \cross; $n=5$: $3125<3125$ \cross; $n=6$: $7776<15625$ \tick.
\item \emph{Claim:} $n^5<5^n$ for all $n\ge6$. We proceed by induction on $n\ge6$.
\item \emph{(Base case.)} $6^5=7776<15625=5^6$.
\item \emph{(Induction step.)} Let $n\ge6$ and suppose $n^5<5^n$. We must prove $(n+1)^5<5^{n+1}$. Since $n\ge6$, $\frac{n+1}n=1+\frac1n\le\frac76$, so $(n+1)^5\le\big(\frac76\big)^5n^5=\frac{16807}{7776}n^5<5n^5$ \why{$\frac{16807}{7776}<5$ since $16807<38880$}. Then $5n^5<5\cdot5^n=5^{n+1}$ \why{induction hypothesis}. So $(n+1)^5<5^{n+1}$.
\item By induction the claim holds for all $n\ge6$. Together with Step 1 this covers every $n\in\N$.
\end{steps}
\answer $n^5<5^n$ exactly for $n\in\{0,1\}$ and for all $n\ge6$.\qed
\begin{compare}
\item The induction only covers $n\ge6$; the values $n\le5$ must be checked one by one and reported.
\item The key estimate is $\big(1+\frac1n\big)^5<5$ for $n\ge6$; justify it with a number, not with ``clearly''.
\end{compare}
\end{bookex}

\begin{bookex}[essential]{4.E.11}
Prove that $(1+x)^{123\,456\,789}\ge1+123\,456\,789\,x$ for all real $x\ge-1$.
\solution
\plan The number is irrelevant (the hint): prove $(1+x)^n\ge1+nx$ for all $n\in\N$ by induction on $n$, for a fixed $x\ge-1$. Induction runs over $n$, never over the real variable $x$.
\begin{steps}
\item Fix $x\in\R$ with $x\ge-1$; then $1+x\ge0$. We prove $p(n)$: $(1+x)^n\ge1+nx$ for all $n\in\N$, by induction on $n$.
\item \emph{(Base case.)} $(1+x)^0=1\ge1=1+0\cdot x$.
\item \emph{(Induction step.)} Let $n\ge0$ and suppose $(1+x)^n\ge1+nx$. Multiply both sides by $1+x\ge0$, which preserves $\ge$ \why{multiplying an inequality by a non-negative number}:
\[
(1+x)^{n+1}=(1+x)(1+x)^n\ge(1+x)(1+nx)=1+nx+x+nx^2=1+(n+1)x+nx^2 .
\]
Since $nx^2\ge0$ \why{$n\ge0$, $x^2\ge0$}, $1+(n+1)x+nx^2\ge1+(n+1)x$. So $(1+x)^{n+1}\ge1+(n+1)x$.
\item By induction $p(n)$ holds for every $n$; take $n=123\,456\,789$. Since $x\ge-1$ was arbitrary, the statement holds for all such $x$.
\end{steps}
\conclusion $(1+x)^{123456789}\ge1+123456789\,x$ for all $x\ge-1$ (Bernoulli's inequality).\qed
\begin{compare}
\item Say why $x\ge-1$ matters: it makes $1+x\ge0$, so multiplying the hypothesis by $1+x$ keeps the direction of the inequality.
\end{compare}
\end{bookex}

\begin{bookex}[recommended]{4.E.12}
Let $a\in\N$ be a number whose decimal expansion ends in the digit $6$. Prove that $a^n$ ends in $6$ for all $n\ge1$.
\solution
\given ``Ends in $6$'' means the last decimal digit is $6$, i.e.\ the remainder on division by $10$ is $6$ (Proposition 1.2.12's idea): $a=10q+6$ for some $q\in\N$.
\begin{steps}
\item Write $a=10q+6$ with $q\in\N$. We proceed by induction on $n\ge1$ with $p(n)$: $a^n=10m+6$ for some $m\in\N$.
\item \emph{(Base case.)} $a^1=a=10q+6$.
\item \emph{(Induction step.)} Let $n\ge1$ and suppose $a^n=10m+6$ with $m\in\N$. Then
\begin{align*}
a^{n+1}=a^n\cdot a&=(10m+6)(10q+6)=100mq+60m+60q+36\\
&=10(10mq+6m+6q+3)+6,
\end{align*}
with $10mq+6m+6q+3\in\N$. So $a^{n+1}$ ends in $6$.
\end{steps}
\conclusion By induction, $a^n$ ends in $6$ for all $n\ge1$.\qed
\begin{compare}
\item Translate ``ends in $6$'' into an equation before inducting; the computation $36=30+6$ is the whole step.
\end{compare}
\end{bookex}

\begin{bookex}[essential]{4.E.13}
Let $a,b\in\R$ and define $a_0=a$, $a_1=b$ and $a_n=\dfrac{a_{n-1}+a_{n+1}}2$ for all $n\ge1$. Prove that $a_n=a+(b-a)n$ for all $n\in\N$.
\solution
\plan The relation defines $a_{n+1}$ from the two previous terms: $a_{n+1}=2a_n-a_{n-1}$ for $n\ge1$. Strong induction, two base cases, step fixed at $n\ge1$.
\begin{steps}
\item Rearranging the defining relation: $a_{n+1}=2a_n-a_{n-1}$ for all $n\ge1$.
\item We proceed by strong induction on $n\ge0$ with $p(n)$: $a_n=a+(b-a)n$, base cases $n=0,1$.
\item \emph{(Base cases.)} $a_0=a=a+(b-a)\cdot0$; $a_1=b=a+(b-a)\cdot1$.
\item \emph{(Induction step.)} Fix $n\ge1$ and assume $a_k=a+(b-a)k$ for all $0\le k\le n$. Then
\begin{align*}
a_{n+1}=2a_n-a_{n-1}&=2\big(a+(b-a)n\big)-\big(a+(b-a)(n-1)\big)\\
&=a+(b-a)(2n-n+1)=a+(b-a)(n+1),
\end{align*}
using the hypothesis for $k=n$ and $k=n-1$.
\end{steps}
\conclusion By strong induction, $a_n=a+(b-a)n$ for all $n\in\N$ (an arithmetic progression).\qed
\end{bookex}

\begin{bookex}[recommended]{4.E.14}
Let $f:\R\to\R$ satisfy $f(x+y)=f(x)+f(y)$ for all $x,y\in\R$. (a) Prove by induction that there is a real number $a$ such that $f(n)=an$ for all $n\in\N$. (b) Deduce $f(n)=an$ for all $n\in\Z$. (c) Deduce $f(x)=ax$ for all $x\in\Q$.
\solution
\begin{steps}
\item \emph{(a)} Define $a=f(1)$. First $f(0)=f(0+0)=f(0)+f(0)$, so $f(0)=0$ \why{subtract $f(0)$}.
\item Induction on $n\in\N$ with $p(n)$: $f(n)=an$. \emph{Base:} $f(0)=0=a\cdot0$. \emph{Step:} if $f(n)=an$ then $f(n+1)=f(n)+f(1)=an+a=a(n+1)$ \why{functional equation; hypothesis}. So $f(n)=an$ for all $n\in\N$.
\item \emph{(b)} Let $n\in\N$. $0=f(0)=f(n+(-n))=f(n)+f(-n)=an+f(-n)$, so $f(-n)=-an=a\cdot(-n)$. Every integer is $n$ or $-n$ for some $n\in\N$, so $f(m)=am$ for all $m\in\Z$.
\item \emph{(c)} First, $f(bx)=bf(x)$ for all $b\in\N$ and $x\in\R$, by induction on $b$: $f(0\cdot x)=f(0)=0$; and $f((b+1)x)=f(bx+x)=f(bx)+f(x)=bf(x)+f(x)=(b+1)f(x)$.
\item Let $x\in\Q$; write $x=\frac mb$ with $m\in\Z$ and $b\in\N$, $b\ge1$ \why{Definition 0.15, choosing the denominator positive}. Then $bf(x)=f(bx)=f(m)=am$ \why{Step 4; (b)}, so $f(x)=\frac{am}b=a\cdot\frac mb=ax$ \why{$b\ne0$}.
\end{steps}
\conclusion $f(x)=ax$ for all $x\in\Q$, with $a=f(1)$.\qed
\begin{compare}
\item Each extension ($\N\to\Z\to\Q$) uses the functional equation once more; (c) needs its own small induction on the denominator.
\end{compare}
\end{bookex}

\begin{bookex}[recommended]{4.E.15}
Let $f:\R\to\R$ satisfy $f(0)>0$ and $f(x+y)=f(x)f(y)$ for all $x,y\in\R$. Prove that there is a positive real number $a$ such that $f(x)=a^x$ for all $x\in\Q$.
\solution
\plan Mirror 4.E.14 with products instead of sums.
\begin{steps}
\item $f(0)=f(0+0)=f(0)^2$ and $f(0)>0$, so $f(0)=1$ \why{divide by $f(0)$}.
\item For every $x$: $f(x)f(-x)=f(0)=1$, so $f(x)\ne0$; and $f(x)=f(\frac x2+\frac x2)=f(\frac x2)^2>0$. So $f$ is positive everywhere. Define $a=f(1)>0$.
\item \emph{Naturals:} induction on $n$: $f(0)=1=a^0$; if $f(n)=a^n$ then $f(n+1)=f(n)f(1)=a^{n+1}$.
\item \emph{Integers:} $f(-n)=\frac1{f(n)}=a^{-n}$ \why{Step 2}.
\item \emph{Rationals:} by induction on $b\in\N$, $f(bx)=f(x)^b$ ($f(0)=1$; $f((b+1)x)=f(bx)f(x)$). For $x=\frac mb$ with $m\in\Z$, $b\ge1$: $f(x)^b=f(bx)=f(m)=a^m$. Since $f(x)>0$, $f(x)$ is the unique positive $b$-th root of $a^m$, which is $a^{m/b}=a^x$ \why{uniqueness of positive roots, Example 1.2.36 generalised}.
\end{steps}
\conclusion $f(x)=a^x$ for all $x\in\Q$ with $a=f(1)>0$.\qed
\begin{compare}
\item Positivity of $f$ (Step 2) is what allows taking the $b$-th root uniquely; it must be proved from $f(0)>0$.
\end{compare}
\end{bookex}

\begin{bookex}[essential]{4.E.16}
Let $a,b\in\Z$. Prove that $b-a$ divides $b^n-a^n$ for all $n\in\N$.
\solution
\given Definition 0.36: $p(n)$ is $\exists k\in\Z,\ b^n-a^n=(b-a)k$.
\begin{steps}
\item We proceed by induction on $n\ge0$.
\item \emph{(Base case.)} $b^0-a^0=1-1=0=(b-a)\cdot0$, so $b-a\mid b^0-a^0$ \why{$k=0$}.
\item \emph{(Induction step.)} Let $n\ge0$ and suppose $b^n-a^n=(b-a)k$ with $k\in\Z$. Then
\begin{align*}
b^{n+1}-a^{n+1}&=b\cdot b^n-a\cdot a^n=b\cdot b^n-b\cdot a^n+b\cdot a^n-a\cdot a^n\\
&=b(b^n-a^n)+(b-a)a^n=(b-a)(bk+a^n),
\end{align*}
\why{add and subtract $b\,a^n$; then the hypothesis} and $bk+a^n\in\Z$. So $b-a\mid b^{n+1}-a^{n+1}$.
\end{steps}
\conclusion By induction, $b-a\mid b^n-a^n$ for all $n\in\N$.\qed
\begin{compare}
\item The trick ``add and subtract $ba^n$'' must appear; the new quotient $bk+a^n$ must be exhibited and declared an integer.
\end{compare}
\end{bookex}

\begin{bookex}[essential]{4.E.17}
Prove by induction that $7^n-2\cdot4^n+1$ is divisible by $18$ for all $n\in\N$.
\solution
\given The book's hint: $7^{n+1}-2\cdot4^{n+1}+1=4(7^n-2\cdot4^n+1)+3(7^n-1)$. (Check: the right side is $4\cdot7^n-8\cdot4^n+4+3\cdot7^n-3=7\cdot7^n-2\cdot4^{n+1}+1$.) So the step also needs $18\mid3(7^n-1)$, i.e.\ $6\mid7^n-1$: a lemma.
\begin{steps}
\item \emph{Lemma: $6\mid7^n-1$ for all $n\in\N$.} Induction: $7^0-1=0=6\cdot0$. If $7^n-1=6m$ then $7^{n+1}-1=7\cdot7^n-1=7(7^n-1)+6=6(7m+1)$ \why{$7\cdot7^n-1=7(7^n-1)+7-1$}.
\item \emph{Main induction} on $n\ge0$ with $p(n)$: $18\mid7^n-2\cdot4^n+1$.
\item \emph{(Base case.)} $7^0-2\cdot4^0+1=1-2+1=0=18\cdot0$.
\item \emph{(Induction step.)} Let $n\ge0$ and suppose $7^n-2\cdot4^n+1=18k$ with $k\in\Z$. By the Lemma, $7^n-1=6m$ with $m\in\Z$. Then
\[
7^{n+1}-2\cdot4^{n+1}+1=4(7^n-2\cdot4^n+1)+3(7^n-1)=4\cdot18k+3\cdot6m=18(4k+m),
\]
and $4k+m\in\Z$. So $18\mid7^{n+1}-2\cdot4^{n+1}+1$.
\end{steps}
\conclusion By induction, $18$ divides $7^n-2\cdot4^n+1$ for all $n\in\N$.\qed
\begin{compare}
\item Verify the hint's identity (it is the whole step), and prove the lemma $6\mid7^n-1$ separately: without it the term $3(7^n-1)$ is not visibly a multiple of $18$.
\end{compare}
\end{bookex}
\begin{bookex}[essential]{4.E.18}
Prove that $\dfrac{d^n}{dx^n}(xe^x)=(n+x)e^x$ for all $n\in\N$.
\solution
\given The $0$-th derivative of a function is the function itself; the product rule $(uv)'=u'v+uv'$; $(e^x)'=e^x$.
\begin{steps}
\item We proceed by induction on $n\ge0$ with $p(n)$: $\frac{d^n}{dx^n}(xe^x)=(n+x)e^x$.
\item \emph{(Base case.)} $\frac{d^0}{dx^0}(xe^x)=xe^x=(0+x)e^x$.
\item \emph{(Induction step.)} Let $n\ge0$ and suppose $\frac{d^n}{dx^n}(xe^x)=(n+x)e^x$. Differentiating once more:
\[
\frac{d^{n+1}}{dx^{n+1}}(xe^x)=\frac d{dx}\big((n+x)e^x\big)=1\cdot e^x+(n+x)e^x=(n+1+x)e^x \why{product rule; hypothesis}.
\]
\end{steps}
\conclusion By induction the formula holds for all $n\in\N$.\qed
\begin{compare}
\item ``Differentiating the hypothesis'' is the step; write the product rule out.
\end{compare}
\end{bookex}

\begin{bookex}[essential]{4.E.19}
Find and prove an expression for $\dfrac{d^n}{dx^n}(x^2e^x)$ for all $n\in\N$.
\solution
\plan Compute the first few derivatives, guess a formula, prove it by induction.
\begin{steps}
\item $n=0$: $x^2e^x$. $n=1$: $(x^2+2x)e^x$. $n=2$: $(2x+2)e^x+(x^2+2x)e^x=(x^2+4x+2)e^x$. $n=3$: $(2x+4)e^x+(x^2+4x+2)e^x=(x^2+6x+6)e^x$.
\item Pattern: $(x^2+2nx+c_n)e^x$ with $c_0,c_1,c_2,c_3=0,0,2,6$, i.e.\ $c_n=n(n-1)$. \emph{Claim:} $\frac{d^n}{dx^n}(x^2e^x)=\big(x^2+2nx+n(n-1)\big)e^x$.
\item \emph{(Base case.)} $n=0$: $(x^2+0+0)e^x=x^2e^x$.
\item \emph{(Induction step.)} Suppose the claim for $n$. Then
\begin{align*}
\frac{d^{n+1}}{dx^{n+1}}(x^2e^x)&=\frac d{dx}\Big(\big(x^2+2nx+n(n-1)\big)e^x\Big)\why{hypothesis}\\
&=(2x+2n)e^x+\big(x^2+2nx+n(n-1)\big)e^x\why{product rule}\\
&=\big(x^2+(2n+2)x+n(n-1)+2n\big)e^x=\big(x^2+2(n+1)x+(n+1)n\big)e^x,
\end{align*}
since $n(n-1)+2n=n^2+n=(n+1)n$. This is the claim for $n+1$.
\end{steps}
\conclusion $\dfrac{d^n}{dx^n}(x^2e^x)=\big(x^2+2nx+n(n-1)\big)e^x$ for all $n\in\N$.\qed
\begin{compare}
\item Show the computations that led to the guess; the proof is the induction, the guess is not.
\end{compare}
\end{bookex}

\begin{bookex}[essential]{4.E.20}
Let $f,g$ be functions that can be differentiated as often as needed. Prove that for all $n\in\N$,
$\dfrac{d^n}{dx^n}\big(f(x)g(x)\big)=\displaystyle\sum_{k=0}^n\binom nkf^{(k)}(x)g^{(n-k)}(x)$.
\solution
\given Product rule; Pascal's rule $\binom n{k-1}+\binom nk=\binom{n+1}k$ (Definition 4.1.15); $\binom nn=1=\binom{n+1}{n+1}$, $\binom n0=1=\binom{n+1}0$. The proof has the same shape as that of the binomial theorem (Theorem 4.2.18).
\begin{steps}
\item Induction on $n\ge0$. \emph{(Base case.)} $n=0$: the right side is $\binom00f^{(0)}g^{(0)}=fg$, the $0$-th derivative of $fg$.
\item \emph{(Induction step.)} Suppose the formula for $n$. Differentiate it term by term with the product rule:
\begin{align*}
\frac{d^{n+1}}{dx^{n+1}}(fg)&=\sum_{k=0}^n\binom nk\Big(f^{(k+1)}g^{(n-k)}+f^{(k)}g^{(n+1-k)}\Big)\\
&=\sum_{k=0}^n\binom nkf^{(k+1)}g^{(n-k)}+\sum_{k=0}^n\binom nkf^{(k)}g^{(n+1-k)} .
\end{align*}
\item Reindex the first sum with $j=k+1$ (so $k=j-1$, $j$ from $1$ to $n+1$): $\sum_{j=1}^{n+1}\binom n{j-1}f^{(j)}g^{(n+1-j)}$. Rename $j$ back to $k$.
\item Split off the extreme terms: from the first sum the term $k=n+1$, namely $\binom nnf^{(n+1)}g^{(0)}$; from the second the term $k=0$, namely $\binom n0f^{(0)}g^{(n+1)}$. The remaining terms, for $1\le k\le n$, combine:
\[
\sum_{k=1}^n\Big(\binom n{k-1}+\binom nk\Big)f^{(k)}g^{(n+1-k)}=\sum_{k=1}^n\binom{n+1}kf^{(k)}g^{(n+1-k)} \why{Pascal's rule}.
\]
\item Since $\binom nn=1=\binom{n+1}{n+1}$ and $\binom n0=1=\binom{n+1}0$, the two extreme terms are the $k=n+1$ and $k=0$ terms of $\sum_{k=0}^{n+1}\binom{n+1}kf^{(k)}g^{(n+1-k)}$. Altogether $\frac{d^{n+1}}{dx^{n+1}}(fg)=\sum_{k=0}^{n+1}\binom{n+1}kf^{(k)}g^{(n+1-k)}$.
\end{steps}
\conclusion By induction, the Leibniz rule holds for all $n\in\N$.\qed
\begin{compare}
\item The three moves are: product rule inside the sum, reindex, Pascal. Each must be visible; the reindexing is where errors hide.
\end{compare}
\end{bookex}

\begin{bookex}[recommended]{4.E.21}
Find and prove an expression for $\dfrac{d^n}{dx^n}\log x$ (for $x>0$) for all $n\in\N$.
\solution
\begin{steps}
\item $n=0$: $\log x$. $n=1$: $x^{-1}$. $n=2$: $-x^{-2}$. $n=3$: $2x^{-3}$. $n=4$: $-6x^{-4}$. Pattern for $n\ge1$: sign $(-1)^{n-1}$, coefficient $(n-1)!$, power $x^{-n}$.
\item \emph{Claim:} for all $n\ge1$, $\frac{d^n}{dx^n}\log x=(-1)^{n-1}(n-1)!\,x^{-n}$. Induction on $n\ge1$.
\item \emph{(Base case.)} $n=1$: $(-1)^0\cdot0!\cdot x^{-1}=\frac1x=\frac d{dx}\log x$.
\item \emph{(Induction step.)} Let $n\ge1$ and suppose the claim for $n$. Then
\[
\frac{d^{n+1}}{dx^{n+1}}\log x=\frac d{dx}\big((-1)^{n-1}(n-1)!\,x^{-n}\big)=(-1)^{n-1}(n-1)!\cdot(-n)x^{-n-1}=(-1)^nn!\,x^{-(n+1)},
\]
which is the claim for $n+1$ \why{$(n-1)!\cdot n=n!$; $(-1)^{n-1}\cdot(-1)=(-1)^n$}.
\end{steps}
\conclusion For $n\ge1$, $\frac{d^n}{dx^n}\log x=(-1)^{n-1}(n-1)!\,x^{-n}$; for $n=0$ it is $\log x$.\qed
\end{bookex}

\begin{bookex}[recommended]{4.E.24}
True or false: every factorial is positive.
\solution
\begin{steps}
\item \emph{True.} Induction on $n$ with $p(n)$: $n!>0$. \emph{Base:} $0!=1>0$ \why{Definition 4.1.14}. \emph{Step:} if $n!>0$ then $(n+1)!=(n+1)\cdot n!>0$ as a product of the positive numbers $n+1\ge1$ and $n!$.
\end{steps}
\conclusion $n!>0$ for all $n\in\N$.\qed
\end{bookex}

\begin{bookex}[essential]{4.E.25}
True or false: every binomial coefficient is positive.
\solution
\begin{steps}
\item \emph{False.} $\binom01=0$ by the clause $\binom0{k+1}=0$ of Definition 4.1.15 (with $k=0$). In general $\binom nk=0$ whenever $k>n$ \why{Theorem 4.2.15}.
\end{steps}
\conclusion False (binomial coefficients are non-negative, not positive).\qed
\end{bookex}

\begin{bookex}[essential]{4.E.26}
True or false: every proof by strong induction can be turned into a proof by weak induction by changing only the induction hypothesis.
\solution
\begin{steps}
\item \emph{False.} A strong-induction step may use $p(k)$ for $k<n$: Example 4.3.7 uses $a_{n-1}$, Example 4.3.4 uses all of $b_0,\dots,b_n$. If the hypothesis is replaced by just $p(n)$, those uses are no longer justified, so the step no longer proves $p(n+1)$: the rest of the argument must change too (for instance by proving the stronger statement $q(n)$: ``$p(k)$ for all $k\le n$'', as in the proof of Theorem 4.3.2).
\end{steps}
\conclusion False. (Every statement provable by strong induction is provable by weak induction, but not by the same proof with only the hypothesis changed.)\qed
\end{bookex}

\begin{bookex}[essential]{4.E.27}
True or false: every proof by weak induction can be turned into a proof by strong induction by changing only the induction hypothesis.
\solution
\begin{steps}
\item \emph{True.} Take a proof by weak induction of $\forall n\ge n_0,\ p(n)$. Replace the hypothesis ``$p(n)$'' by ``$p(k)$ for all $n_0\le k\le n$''.
\item The new hypothesis contains the old one (take $k=n$), so every line of the original step remains valid and still derives $p(n+1)$; the base case is unchanged. This is a valid proof by strong induction \why{Strategy 4.3.3}.
\end{steps}
\conclusion True.\qed
\end{bookex}

\begin{bookex}[recommended]{4.E.28}
True or false: let $f:\N\to X$ be a surjection and $p(x)$ a logical formula with free variable $x\in X$. If $p(f(0))$ is true and $\forall n\in\N,\ (p(f(n))\Rightarrow p(f(n+1)))$, then $\forall x\in X,\ p(x)$.
\solution
\begin{steps}
\item \emph{True.} Define $q(n)$: $p(f(n))$, a formula with free variable $n\in\N$. The hypotheses say $q(0)$ and $\forall n,\ (q(n)\Rightarrow q(n+1))$.
\item By weak induction, $q(n)$ holds for all $n\in\N$ \why{Theorem 4.2.1}.
\item Let $x\in X$. Since $f$ is surjective, $x=f(n)$ for some $n\in\N$ \why{Definition 3.2.9}. Then $p(x)$ is $p(f(n))=q(n)$, which is true.
\end{steps}
\conclusion True: induction transports along a surjection.\qed
\end{bookex}

\begin{bookex}[essential]{4.E.29}
True or false: if for all $x\in\N$ there is some $y\le x$ such that $p(y)$ is false, then $p(x)$ is false for all $x\in\N$.
\solution
\begin{steps}
\item \emph{False.} Counterexample: $p(x)$: ``$x\ge1$''.
\item \emph{Hypothesis holds:} for every $x\in\N$, $y=0$ satisfies $y\le x$ and $p(0)$ (``$0\ge1$'') is false.
\item \emph{Conclusion fails:} $p(1)$ is true.
\end{steps}
\conclusion False. (The hypothesis only guarantees that $p(0)$ is false; it is not an induction principle.)\qed
\end{bookex}

\begin{bookex}[essential]{4.E.30}
True or false: every non-empty subset of $\Q_{\ge0}$ has a least element.
\solution
\begin{steps}
\item \emph{False.} Let $X=\{\frac1n\mid n\in\N,\ n\ge1\}\subseteq\Q_{\ge0}$; it is non-empty ($1\in X$).
\item $X$ has infinite descent: for $\frac1n\in X$, the element $\frac1{n+1}\in X$ satisfies $\frac1{n+1}<\frac1n$ \why{Definition 4.3.16}. A set with infinite descent has no least element (a least element $m$ would have some $y\in X$ with $y<m$).
\end{steps}
\conclusion False: well-ordering is special to $\N$.\qed
\end{bookex}

\begin{bookex}[essential]{4.E.31}
True or false: every non-empty subset of $\Z$ has a least element.
\solution
\begin{steps}
\item \emph{False.} $\Z$ itself is a non-empty subset of $\Z$, and it has infinite descent: for $n\in\Z$, $n-1\in\Z$ and $n-1<n$ \why{Example 4.3.17}. So $\Z$ has no least element.
\end{steps}
\conclusion False.\qed
\end{bookex}

\begin{bookex}[essential]{4.E.32}
True or false: every non-empty subset of $\Z^-$ (the set of negative integers) has a greatest element.
\solution
\plan Flip the set into $\N$ and use the well-ordering principle.
\begin{steps}
\item \emph{True.} Let $X\subseteq\Z^-$ be non-empty. Define $Y=\{-x-1\mid x\in X\}$. For $x\le-1$ we have $-x-1\ge0$, so $Y\subseteq\N$; and $Y\ne\varnothing$ since $X\ne\varnothing$.
\item By the well-ordering principle \why{Theorem 4.3.11}, $Y$ has a least element, say $-x_0-1$ with $x_0\in X$.
\item Let $x\in X$. Then $-x-1\in Y$, so $-x-1\ge-x_0-1$, hence $x\le x_0$ \why{multiply by $-1$}. So $x_0$ is the greatest element of $X$.
\end{steps}
\conclusion True.\qed
\begin{compare}
\item The map $x\mapsto-x-1$ (or $x\mapsto-x$, landing in $\N\setminus\{0\}$) must be shown to land in $\N$, and ``least'' must be translated back to ``greatest''.
\end{compare}
\end{bookex}

\begin{bookex}[recommended]{4.E.33}
Let $n,k,\ell\in\N$. Then $\binom n{k+\ell}=\binom nk+\binom n\ell$. (Always, sometimes, never? The book writes $n,k,\ell\in\Z$; binomial coefficients are defined for natural numbers, so we take $\N$.)
\solution
\begin{steps}
\item \emph{Holds:} $n=0$, $k=\ell=1$: $\binom02=0$ and $\binom01+\binom01=0+0=0$ \why{Theorem 4.2.15}.
\item \emph{Fails:} $n=1$, $k=0$, $\ell=1$: $\binom11=1$ but $\binom10+\binom11=1+1=2$.
\end{steps}
\conclusion Sometimes.\qed
\end{bookex}

\begin{bookex}[recommended]{4.E.34}
Let $m,n\ge2$. Then $(m+n)!=m!+n!$. (Always, sometimes, never?)
\solution
\plan Show $(m+n)!>m!+n!$ always: $(m+n)!$ contains both $m!$ and $n!$ as factors of a much larger product.
\begin{steps}
\item \emph{Never.} First, if $j\le\ell$ then $j!\le\ell!$ \why{$\ell!=j!\cdot(j+1)\cdots\ell$ and each extra factor is $\ge1$; induction on $\ell-j$}.
\item Since $m+n-1\ge m$ and $m+n-1\ge n$ \why{$m,n\ge2$, so $n-1\ge1$ and $m-1\ge1$}, Step 1 gives $(m+n-1)!\ge m!$ and $(m+n-1)!\ge n!$, so $(m+n-1)!\ge\max(m!,n!)$.
\item Then $(m+n)!=(m+n)(m+n-1)!\ge4\max(m!,n!)>2\max(m!,n!)\ge m!+n!$ \why{$m+n\ge4$; $\max(m!,n!)>0$ by 4.E.24}.
\end{steps}
\conclusion $(m+n)!>m!+n!$ for all $m,n\ge2$: never.\qed
\begin{compare}
\item ``Never'' needs an argument valid for \emph{all} $m,n\ge2$, not a few numerical checks.
\end{compare}
\end{bookex}

\begin{bookex}[recommended]{4.E.35}
Let $p(x)$ be a logical formula with free variable $x\in\Z$. Suppose $p(0)$ is true and, for all $n\in\Z$, $p(n)\Rightarrow p(n+1)$ and $p(n)\Rightarrow p(n-2)$. Then $\forall n\in\Z,\ p(n)$. (Always, sometimes, never?)
\solution
\plan Non-negative integers by weak induction; negative ones by strong induction on $m$ for $q(m)$: $p(-m)$, using $p(n)\Rightarrow p(n-2)$ to go from $-(m-1)$ to $-(m+1)$.
\begin{steps}
\item \emph{Always.} \emph{Non-negative:} $p(0)$ holds and $p(n)\Rightarrow p(n+1)$ for all $n$, so $p(n)$ for all $n\in\N$ by weak induction \why{Theorem 4.2.1}.
\item \emph{Negative:} let $q(m)$ be $p(-m)$ for $m\in\N$. Strong induction with base cases $m=0,1$: $q(0)=p(0)$ holds; $q(1)=p(-1)$ follows from $p(1)$ (Step 1) via $p(1)\Rightarrow p(1-2)$.
\item \emph{Step:} fix $m\ge1$ and assume $q(k)$ for all $0\le k\le m$. Then $q(m-1)=p(-(m-1))$ holds, and $p(-(m-1))\Rightarrow p(-(m-1)-2)=p(-(m+1))=q(m+1)$. So $q(m+1)$.
\item Hence $p(-m)$ for all $m\in\N$. Every integer is $n\in\N$ or $-m$ with $m\in\N$, so $p$ holds on all of $\Z$.
\end{steps}
\conclusion Always.\qed
\end{bookex}

